% Source: https://artofproblemsolving.com/community/c3h3167127_2022_lmt_spring_guts_round_p1p15_lexington_mathematical_tournament
% Stable identity: aops-topic-3167127-q13
Triangle $ABD$ is right, so $BD=\dfrac{4\sqrt3}{\cos25^\circ}$ and its area is $\dfrac{1}{2}AB^2\tan25^\circ=24\tan25^\circ$. In triangle $BCD$, the remaining angle at $B$ is $115^\circ$. By the sine rule, $CD=BD\dfrac{\sin115^\circ}{\sin60^\circ}=8$. Its area is $\dfrac{1}{2}BD\cdot8\sin5^\circ=\dfrac{16\sqrt3\sin5^\circ}{\cos25^\circ}$. Since $\sin5^\circ=\sin(30^\circ-25^\circ)=\dfrac{1}{2}\cos25^\circ-\dfrac{\sqrt3}{2}\sin25^\circ$, this second area is $8\sqrt3-24\tan25^\circ$. The sum is $\boxed{8\sqrt3}$.
\begin{asy}
usepackage("amsmath");
size(7cm);
real L=4*sqrt(3)/cos(25*pi/180);pair Bp=(0,0),D=(L,0),A=4*sqrt(3)*dir(25),C=D+8*dir(185);draw(A--Bp--C--D--cycle);draw(Bp--D,dashed);label("$A$",A,N);label("$B$",Bp,W);label("$C$",C,S);label("$D$",D,E);
\end{asy}
